\section{Check the laptop}
\label{sec:prelim}

These three checks take five minutes and prevent the most expensive kind of failure, an
installation that cannot work on the machine at all.

\step{Find the Windows version}
Hold the \key{Windows} key (the key with the four-pane logo, next to \key{Alt}) and press
\key{R}. A small window titled \ui{Run} opens. Type \code|winver| in its box and press
\key{Enter}. A window titled \ui{About Windows} appears; read its first lines and then close it
with \ui{OK}.
\checkpoint{The window says \emph{Windows 10} or \emph{Windows 11}. For Windows 10 the version
shown must be 1809 or later (for example 22H2).}

\step{Find the processor type}
Click the \ui{Start} button, type \code|about your pc| and press \key{Enter}. The Settings
window opens on a page headed \ui{About}. Under \ui{Device specifications}, find the line
\ui{System type}.
\checkpoint{The line reads \emph{64-bit operating system, x64-based processor}.}

\begin{iffail}[If the line says ``ARM-based processor'' or ``32-bit'']
Stop here and tell the project lead. The neural-network library used by the project
(TensorFlow) is published for x64 processors only. The rest of the software can still be
installed, and the project's simpler fallback classifier works without TensorFlow, but the lead
should decide how your laptop is used.
\end{iffail}

\step{Find the free disk space}
Hold \key{Windows} and press \key{E} to open File Explorer. In its left-hand list click
\ui{This PC}. Under \ui{Devices and drives}, read the text below the drive named
\ui{Local Disk (C:)} or \ui{Windows (C:)}.
\checkpoint{At least 10\,GB is free on drive C:.}

\step{Show file name extensions}
The guide refers to files by their full names, such as \file{config.py} or \file{T26.png}.
Windows hides the part after the dot unless told otherwise. In File Explorer on Windows~11,
click \menu{View}\then\menu{Show}\then\menu{File name extensions}; on Windows~10, click the
\menu{View} tab and tick the box \ui{File name extensions}.
\checkpoint{File names in File Explorer now end in a dot and a few letters, for example
\file{.pdf} or \file{.png}.}
